\documentclass[11pt]{article}
\usepackage[T1]{fontenc}
\usepackage[utf8]{inputenc}
\usepackage{amsmath,amsthm,mathtools}
\usepackage{newtxtext,newtxmath}
\usepackage[margin=1in,headheight=15pt]{geometry}
\usepackage{microtype}
\usepackage{xcolor}
\usepackage{booktabs,array}
\usepackage{enumitem}
\usepackage{fancyhdr}
\usepackage{hyperref}
\usepackage[nameinlink,noabbrev]{cleveref}
\definecolor{inkblue}{HTML}{173A5E}
\definecolor{linkblue}{HTML}{205C83}
\hypersetup{colorlinks=true,linkcolor=linkblue,citecolor=linkblue,urlcolor=linkblue,
 pdftitle={Growth-Scale Rigidity for Surreal and Surcomplex Automorphisms},
 pdfauthor={Research manuscript prepared with ChatGPT},
 pdfsubject={Exponential automorphisms, convex valuations, and a negative answer to Question 5.4 of Kaplan--Krapp--Serra}}
\pagestyle{fancy}
\fancyhf{}
\fancyhead[L]{\small\scshape Growth-scale rigidity}
\fancyhead[R]{\small Surreal and surcomplex automorphisms}
\fancyfoot[C]{\thepage}
\renewcommand{\headrulewidth}{0.35pt}
\setlength{\parskip}{3pt}
\setlength{\parindent}{1.2em}
\setlist{nosep,leftmargin=1.6em}
\newtheorem{theorem}{Theorem}[section]
\newtheorem{lemma}[theorem]{Lemma}
\newtheorem{proposition}[theorem]{Proposition}
\newtheorem{corollary}[theorem]{Corollary}
\theoremstyle{definition}
\newtheorem{definition}[theorem]{Definition}
\newtheorem{example}[theorem]{Example}
\theoremstyle{remark}
\newtheorem{remark}[theorem]{Remark}
\newtheorem{question}[theorem]{Question}
\crefname{theorem}{Theorem}{Theorems}
\crefname{lemma}{Lemma}{Lemmas}
\crefname{proposition}{Proposition}{Propositions}
\crefname{corollary}{Corollary}{Corollaries}
\crefname{question}{Question}{Questions}
\crefname{definition}{Definition}{Definitions}
\newcommand{\No}{\mathbf{No}}
\newcommand{\R}{\mathbb{R}}
\newcommand{\Q}{\mathbb{Q}}
\newcommand{\N}{\mathbb{N}}
\newcommand{\C}{\mathbb{C}}
\newcommand{\OO}{\mathcal{O}}
\newcommand{\mm}{\mathfrak{m}}
\newcommand{\id}{\operatorname{id}}
\newcommand{\Aut}{\operatorname{Aut}}
\newcommand{\supp}{\operatorname{supp}}
\newcommand{\im}{\operatorname{im}}
\newcommand{\st}{\operatorname{st}}
\newcommand{\lt}{\operatorname{lt}}
\newcommand{\cis}{\operatorname{cis}}
\newcommand{\abs}[1]{\lvert #1\rvert}
\newcommand{\norm}[1]{\lVert #1\rVert}
\newcommand{\Gexp}{\operatorname{Aut}_{\exp}}
\newcommand{\GLmod}{\operatorname{Aut}_{L}}
\newcommand{\eqdef}{\coloneqq}
\numberwithin{equation}{section}
\title{\color{inkblue}\bfseries Growth-Scale Rigidity for Surreal\texorpdfstring{\\}{ }and Surcomplex Automorphisms\\[0.5em]
\large A negative answer to Kaplan--Krapp--Serra's Question 5.4,\\
with convex-valuation and non-lifting theorems}
\author{Research manuscript prepared with ChatGPT}
\date{21 September 2026}
\begin{document}
\maketitle
\thispagestyle{empty}
\begin{abstract}
Let $F$ be an ordered field with a total increasing exponential isomorphism
$E:(F,+)\to(F_{>0},\cdot)$, and let $w$ be any nontrivial convex valuation.
We prove that exponential automorphisms stabilizing the valuation ring act
faithfully on the value group. More strongly, every nonidentity such
automorphism induces a value-group automorphism whose additive displacement
is cofinal and coinitial. The same conclusion holds after any nontrivial
invariant convex coarsening. Applied to the surreal exponential field, this
rules out all nonidentity exponential $1$-automorphisms, without a strong
additivity or coefficient-fixing hypothesis, and answers Question~5.4 in
Kaplan--Krapp--Serra, arXiv:2509.22374v3, negatively.

We then exhibit extensive families of ordered value-group automorphisms that
have explicit Hahn-field lifts but no exponential lift. These include every
nonidentity reweighting confined to a bounded collection of normal-form layers,
and every nonidentity positive rational dilation. Finally, for
$K=F(i)$ equipped with $L(z)=\log\abs z$, we determine its automorphisms and show
that the kernel of their action on any invariant nontrivial convex value
quotient is exactly complex conjugation together with the identity.
Along the way we record a common generalization of the displacement identity
and the Leibniz rule: no nonzero twisted derivation of a nontrivially valued
field has all of its output values bounded below by one group element, at any
valuation rank and with no ordering. From it we obtain a form of the rigidity
theorem needing only one visible exponential scale, and hence no surjectivity
of $E$; a form whose hypotheses are imposed only on a final interval; and the
cofinality of the multiplicative motion and of the exponential commutation
defects. A differential counterpart forbids a global valuation-gain bound
$v(\partial y)\ge v(y)+\gamma$ for every fixed $\gamma$, and explicit formal
Laurent-series examples delimit the hypotheses.
All rigidity proofs are algebraic and order-theoretic; none exchanges an
automorphism with an infinite sum.
\end{abstract}
\begin{center}
\small\textbf{Status.} Unrefereed research draft. The proofs are supplied in full;
no proof-assistant certification or exhaustive priority claim is made.
\end{center}
\vspace{0.3em}
\noindent\textbf{Keywords.} Surreal numbers; exponential fields; field automorphisms;
convex valuations; Hahn series; surcomplex numbers; lifting obstructions.
\clearpage
\begingroup
\small
\setlength{\parskip}{0pt}
\tableofcontents
\endgroup
\clearpage

\section{Introduction and the question being answered}

A field automorphism can preserve the leading growth of every element while
changing lower-order terms. Exponentiation makes this much harder: it can
turn an additive error into a change of multiplicative scale. The purpose of
this article is to turn that observation into an exact rigidity theorem,
and to show that it determines much more than the first leading coefficient.

Kaplan, Krapp, and Serra ask in \cite[Question~5.4]{KKS} whether the surreal
exponential field has a nonidentity $1$-automorphism. Their
Proposition~5.2 rules this out under strong $\R$-linearity; the unrestricted
question is printed on page~10 of version~3. A $1$-automorphism fixes every
leading term. We prove the stronger assertion that fixing every valuation
value already forces the identity. The argument does not assume strong
additivity, fixation of the embedded real field, or representation as the
exponential of a derivation.

\subsection{The principal results}

We use additive notation for valuations. Elements outside the valuation ring
have negative values; nonzero elements of its maximal ideal have positive
values.
For an ordered field $F$, an \emph{ordered exponential} will always mean an
increasing group isomorphism
\[
 E:(F,+)\longrightarrow(F_{>0},\cdot).
\]
In particular, $E$ is onto and has an increasing inverse, denoted by $\log$.
If $w:F^\times\to\Gamma$ is a convex valuation and an automorphism $\sigma$
stabilizes its valuation ring, its induced value-group action is denoted by
$\tau_\sigma$, so that
\[
 w(\sigma y)=\tau_\sigma(w(y))\qquad(y\ne0).
\]

\begin{theorem}[Main rigidity theorem]\label[theorem]{thm:intro}
Suppose $F$ has an ordered exponential $E$, $w$ is nontrivial and convex,
and $\sigma$ is an exponential field automorphism stabilizing $\OO_w$.
If $\sigma\ne\id$, the image
\[
 \{\tau_\sigma(\gamma)-\gamma:\gamma\in\Gamma\}
\]
is cofinal and coinitial in $\Gamma$. Consequently, the action on $\Gamma$
is faithful. For every proper convex subgroup $\Delta\subsetneq\Gamma$,
the subgroup stabilizing $\Delta$ acts faithfully on $\Gamma/\Delta$.
\end{theorem}

The coarsening conclusion requires that the automorphisms under consideration
stabilize the chosen convex subgroup. It does not assert that every
exponential automorphism stabilizes every coarsening.

For the surreal field, the natural value group is canonically the ordered
additive group $(\No,+,<)$. Write
\[
 v\left(\sum_\gamma a_\gamma t^\gamma\right)
   =\min\supp(a),\qquad t^\gamma=\omega^{-\gamma}.
\]
The resulting consequences include
\begin{equation}\label{eq:main-no}
 \ker\!\left(\Gexp(\No)\longrightarrow\Aut(\No,+,<)\right)=\{\id\},
 \qquad 1\text{-}\Gexp(\No)=\{\id\}.
\end{equation}
The second equality is the negative answer to the named question. The first
is stronger: it does not require preservation of leading coefficients.

The same theorem yields explicit \emph{non-lifting} statements. At the level
of the value group, independently rescale real coefficients at any bounded-above
collection of normal-form exponents. This gives an ordered additive
automorphism and hence a canonical ordered Hahn-field automorphism. Unless
all rescalings are trivial, however, \emph{no} lift can preserve the surreal
exponential. Positive rational dilations of the whole value group also fail
to lift exponentially, although their displacements are cofinal. Thus
cofinal displacement is necessary but not sufficient for exponential lifting.

Several refinements of the hypotheses are proved along the way and are
independent of the surreal application. \Cref{thm:visible-scale} needs only
one $H>0$ whose exponential is visible to the valuation, in the sense that
$w(E(H))<0$; it therefore does not assume that $E$ is onto $F_{>0}$, and
\cref{cor:detector} states the same mechanism with no ordering and in
arbitrary characteristic. \Cref{thm:profile} replaces both hypotheses of
\cref{thm:no} by equality of the exponential profiles
$w(E(\sigma x))=w(E(x))$ on an additive generating set, and in particular on
one final interval. \Cref{thm:sigma-derivation} isolates the algebra behind
all of this: displacements and derivations are both twisted derivations, and
no nonzero twisted derivation of a nontrivially valued field has its output
values bounded below by a single group element. Beyond rigidity,
\cref{thm:motion} shows that the multiplicative motion of a nonidentity
exponential automorphism is cofinal and coinitial, and \cref{thm:defect}
that the same holds for the valuations of the exponential commutation
defects of a nonidentity automorphism fixing every value.

For the surcomplex field $\No[i]$, define the everywhere-defined, nonzero-input
operation
\[
 L(z)=\log\sqrt{z\bar z}\qquad(z\ne0).
\]
It uses only the real surreal logarithm, not a choice of complex logarithm.
We prove
\begin{equation}\label{eq:main-complex}
 \GLmod(\No[i])\simeq\Gexp(\No)\times C_2,
\end{equation}
and the kernel of the action on a nontrivial invariant convex value quotient
is exactly $\{\id,\text{conjugation}\}$. Marking $i$ removes this final ambiguity.

\subsection{Relation to the repository and novelty scope}

The inspected repository documentation already has a substantial analysis
program, including contour theory, differential equations, and finite spectral
theory \cite{Repo}. We therefore study the automorphism problem rather than
repackage those topics. The reference snapshot is specified in the bibliography;
this is an independent article, not an edit to that repository.

The concrete literature comparison is the explicit Question~5.4 above.
The general cofinal-displacement theorem, its coarsening consequence, the
non-lifting families, and the logarithmic-modulus application are presented
here as further results. Targeted searches did not identify these statements
in this form. This is not a guarantee of priority: the elementary displacement
lemma may be familiar in other language, and a specialist may recognize
additional antecedents. The distinction between proving a statement and
establishing its historical novelty is important.

\section{Conventions and the exponential--valuation bridge}

\subsection{Convex valuations}

Let $F$ be an ordered field. A valuation is a surjective homomorphism
$w:F^\times\to\Gamma$ to an ordered abelian group satisfying
\[
 w(x+y)\ge\min\{w(x),w(y)\},
\]
whenever the expression is defined; set $w(0)=\infty$. Its valuation ring,
maximal ideal, and group of units are
\[
 \OO_w=\{x:w(x)\ge0\},\qquad
 \mm_w=\{x:w(x)>0\},\qquad
 \OO_w^\times=\{x:w(x)=0\}.
\]
Convexity means that $\OO_w$ is order-convex in $F$. Equivalently,
\begin{equation}\label{eq:convex-monotone}
 0<x\le y\quad\Longrightarrow\quad w(x)\ge w(y).
\end{equation}
For completeness, $0<x/y\le1$ belongs to a convex valuation ring, giving one
direction. Conversely, \eqref{eq:convex-monotone} puts every element between
$0$ and an element of $\OO_w$ back in $\OO_w$; signs give convexity.
All nonzero rational numbers are units for a convex valuation. Indeed, for
$n\ge1$, the valuation inequality gives $w(n)\ge0$, while
$1/n\in[0,1]$ gives $w(1/n)\ge0$.

If $\sigma(\OO_w)=\OO_w$, then $\sigma$ also preserves its units.
Consequently the formula
\[
 \tau_\sigma(w(y))=w(\sigma(y))
\]
is well-defined. It is additive and bijective. Since
$w(y)\ge w(z)$ is equivalent to $y/z\in\OO_w$, it preserves the order.
This constructs the induced value-group automorphism without any choice of
monomial representatives.

We will distinguish carefully between \emph{stabilizing the valuation ring}
and \emph{fixing the values pointwise}. The former allows
$\tau_\sigma\ne\id$; the latter says $w(\sigma y)=w(y)$ for every $y\ne0$.
Faithfulness concerns the kernel of this action, not the triviality of the
whole automorphism group.

\subsection{The logarithmic unit subgroup}

\begin{lemma}[The bridge]\label[lemma]{lem:bridge}
Let $E$ be an ordered exponential on $F$, and let $w$ be a nontrivial convex
valuation. The map
\begin{equation}\label{eq:nu}
 \nu:F\longrightarrow\Gamma,\qquad \nu(x)=w(E(x)),
\end{equation}
is a surjective additive, order-reversing map. Its kernel
\begin{equation}\label{eq:H}
 H_w=\{x:E(x)\in\OO_w^\times\}
\end{equation}
is a proper convex additive subgroup of $F$.
\end{lemma}
\begin{proof}
The exponential law gives
$\nu(x+y)=w(E(x)E(y))=\nu(x)+\nu(y)$.
If $x\le y$, then $E(x)\le E(y)$, so convexity of $w$ gives
$\nu(x)\ge\nu(y)$.
For $\gamma\in\Gamma$, choose $a\ne0$ with $w(a)=\gamma$.
Then $\nu(\log\abs a)=\gamma$. This proves surjectivity.
The kernel is an additive subgroup. If $h_1\le x\le h_2$ and
$\nu(h_1)=\nu(h_2)=0$, monotonicity forces $\nu(x)=0$.
Nontriviality of $\Gamma$ and surjectivity imply that the kernel is proper.
\end{proof}

\begin{lemma}[Proper convex subgroups are bounded]\label[lemma]{lem:bounded}
If $H$ is a proper convex additive subgroup of an ordered field $F$, there
exists $B>0$ such that $\abs h<B$ for every $h\in H$.
\end{lemma}
\begin{proof}
Choose $B>0$ outside $H$. If some $\abs h\ge B$, convexity between $0$ and
$\abs h\in H$ would put $B$ in $H$.
\end{proof}

Here ``bounded'' means bounded by an element of $F$, not necessarily by an
integer or a real number. For example, all finite surreal numbers are bounded
by $\omega$. This elementary distinction is decisive.

\subsection{The natural valuation on the surreals}

The standard facts needed about surreal numbers are their real-closed ordered
field structure, Conway normal forms, and the total increasing exponential
extending the real exponential \cite{Gonshor,DriesEhrlich}. In normal form,
\begin{equation}\label{eq:normal}
 x=\sum_{s\in S}r_s\omega^s,
\end{equation}
where $S$ is a reverse well-ordered \emph{set} and $r_s\in\R^\times$.
For $x\ne0$,
\[
 v(x)=-\max S.
\]
The natural valuation ring is
\begin{equation}\label{eq:finite-ring}
 \OO=\{x\in\No:\abs x<n\text{ for some }n\in\N_{>0}\}.
\end{equation}
A positive unit $u$ is characterized by bounds
$1/n<u<n$ for a sufficiently large positive integer $n$.

\begin{proposition}\label[proposition]{prop:finite-log}
For the surreal exponential and natural valuation,
\[
 \ker(v\circ\exp)=\OO.
\]
\end{proposition}
\begin{proof}
If $-n<x<n$, monotonicity gives
$e^{-n}<\exp x<e^n$. Since the real numbers $e^n$ are finite, $\exp x$
and its reciprocal are finite, hence $v(\exp x)=0$.
Conversely, if $\exp x$ is a positive unit, choose an integer $n>1$ with
$n^{-1}<\exp x<n$. Taking logarithms gives
$-\log n<x<\log n$, so $x$ is finite.
\end{proof}

No normal-form formula for $\exp$ is required in this proposition. In
particular, the proof never replaces the Gonshor exponential by the
omega-map. The latter has different algebraic and asymptotic properties.

\section{An algebraic displacement-amplification lemma}

\subsection{The exact identity}

Let $\sigma:F\to F$ be a unital field endomorphism, and put
\[
 D_\sigma(x)=\sigma(x)-x.
\]
The map $D_\sigma$ is additive, but it is generally \emph{not} a derivation.
The correct product identity is
\begin{equation}\label{eq:twisted}
 D_\sigma(ab)=\sigma(a)D_\sigma(b)+bD_\sigma(a).
\end{equation}
Indeed, the right-hand side expands to
$\sigma(a)\sigma(b)-\sigma(a)b+b\sigma(a)-ba$.
Keeping $\sigma(a)$ rather than replacing it by $a$ is essential.

\begin{lemma}[Quantitative displacement amplification]\label[lemma]{lem:amplify}
Let $F$ be an ordered field and $\sigma:F\to F$ a unital field endomorphism.
Suppose $a\in F$ is moved by $\sigma$, and set $d=D_\sigma(a)\ne0$.
For every $B>0$, define
\begin{equation}\label{eq:b}
 b=\frac{2B(1+\abs{\sigma(a)})}{\abs d}.
\end{equation}
Then
\begin{equation}\label{eq:amplify}
 \max\{\abs{D_\sigma(b)},\abs{D_\sigma(ab)}\}>B.
\end{equation}
Consequently, if $\sigma\ne\id$, its additive displacement image
$D_\sigma(F)$ is cofinal and coinitial in $F$.
\end{lemma}
\begin{proof}
If both absolute values in \eqref{eq:amplify} were at most $B$,
\eqref{eq:twisted} and the triangle inequality would give
\begin{align*}
 2B(1+\abs{\sigma(a)})
 &=\abs{bd}\\
 &=\abs{D_\sigma(ab)-\sigma(a)D_\sigma(b)}\\
 &\le B+\abs{\sigma(a)}B,
\end{align*}
a contradiction in an ordered field.
Thus absolute displacements have no upper bound in $F$.
Since $D_\sigma(-x)=-D_\sigma(x)$, both signs occur at every required
magnitude. Given $c\in F$, apply the result with $B>\abs c$ to obtain a
displacement greater than $c$ and another smaller than $c$.
\end{proof}

Surjectivity and order preservation of $\sigma$ are not needed in this lemma.
There are only two probe arguments, $b$ and $ab$, and all expressions are
finite field expressions. In particular, the proof works for self-embeddings
and for class-sized fields without a transfinite construction.

\begin{corollary}[Bounded-displacement rigidity]\label[corollary]{cor:bounded-displacement}
A unital field endomorphism of an ordered field whose displacement image is
contained in a proper convex additive subgroup is the identity.
More generally, a uniform field-valued bound on all absolute displacements
forces the identity.
\end{corollary}
\begin{proof}
Use \cref{lem:bounded} and \cref{lem:amplify}.
\end{proof}

\subsection{A valued-field form of the same principle}

The ordered-field proof is especially transparent. A second formulation
shows that the underlying effect is not specifically tied to an ordering.

\begin{proposition}\label[proposition]{prop:valued-amplification}
Let $(F,w)$ be a nontrivially valued field and let $\sigma:F\to F$ be a
nonidentity field endomorphism. The set
\[
 \{w(\sigma(x)-x):\sigma(x)\ne x\}
\]
is coinitial in the value group.
\end{proposition}
\begin{proof}
Choose $a$ with $d=\sigma(a)-a\ne0$; then $a\ne0$ and $\sigma(a)\ne0$.
Given $\gamma$ in the value group, choose $b\ne0$ so that
\[
 w(b)+w(d)<\gamma+\min\{0,w(\sigma(a))\}.
\]
Such a $b$ exists because a nonzero ordered abelian group has no least element
and $w$ is onto. If both $w(D_\sigma(b))$ and $w(D_\sigma(ab))$ were at least
$\gamma$, the valuation inequality applied to
$bd=D_\sigma(ab)-\sigma(a)D_\sigma(b)$ would give the reverse weak inequality.
Thus one of these two displacement values is less than $\gamma$.
\end{proof}

The proposition neither assumes nor concludes that $\sigma$ preserves the
valuation. It concerns the size of the additive error, not the induced
value-group action. Exponentiation connects those two different notions.

\subsection{One identity for displacement and derivation}

The twisted identity \eqref{eq:twisted} and the ordinary Leibniz rule are
two instances of a single law, and both forms of the amplification argument
use only that law. Recording the common generalization costs nothing, and it
makes the same two probes serve the differential statements of
\cref{sec:differential} as well as the displacement statements above.

\begin{definition}[$\sigma$-derivation]\label[definition]{def:sigma-derivation}
Let $\sigma:F\to F$ be a unital field endomorphism. An additive map
$\delta:F\to F$ is a \emph{$\sigma$-derivation} if
\begin{equation}\label{eq:sigma-leibniz}
 \delta(ab)=\sigma(a)\delta(b)+b\delta(a)\qquad(a,b\in F).
\end{equation}
\end{definition}

Taking $\delta=D_\sigma$ recovers \eqref{eq:twisted}; taking $\sigma=\id$
recovers the Leibniz rule, so ordinary derivations are exactly the
$\id$-derivations. The twist sits on the first factor only, and it is the
same twist as in \eqref{eq:twisted}.

\begin{theorem}[No globally valuation-bounded nonzero $\sigma$-derivation]
\label[theorem]{thm:sigma-derivation}
Let $F$ be a field of arbitrary characteristic carrying a nontrivial
surjective valuation $v:F^\times\to\Gamma$, let $\sigma:F\to F$ be a unital
field endomorphism, and let $\delta$ be a $\sigma$-derivation. If there is a
single $\beta\in\Gamma$ with
\begin{equation}\label{eq:uniform-output}
 v(\delta(x))\ge\beta\qquad(x\in F),
\end{equation}
then $\delta=0$. Equivalently, if $\delta\ne0$, then for every
$\beta\in\Gamma$ some $x$ satisfies $v(\delta(x))<\beta$.
Neither an ordering nor convexity is assumed, no restriction is placed on
the rank of $\Gamma$, and $\sigma$ is not assumed to preserve $v$.
\end{theorem}
\begin{proof}
Assume \eqref{eq:uniform-output} and choose $a$ with $\delta(a)\ne0$.
Additivity gives $\delta(0)=0$, so $a\ne0$; a unital field endomorphism is
injective, so $\sigma(a)\ne0$. Put
\[
 M=\min\{\beta,\;v(\sigma(a))+\beta\}.
\]
For every $b$, rearranging \eqref{eq:sigma-leibniz} and applying the
valuation inequality gives
\[
 v(b\delta(a))=v\bigl(\delta(ab)-\sigma(a)\delta(b)\bigr)
 \ge\min\{v(\delta(ab)),\,v(\sigma(a))+v(\delta(b))\}\ge M.
\]
Choose $\eta>0$ in $\Gamma$, which is possible because $\Gamma\ne0$, and use
surjectivity of $v$ to choose $b\ne0$ with $v(b)=M-v(\delta(a))-\eta$. Then
$v(b\delta(a))=M-\eta<M$, a contradiction.
\end{proof}

The valuation inequality is used only in the safe direction, bounding the
value of a difference below by the minimum. Nothing is assumed about
whether that minimum is attained, so no cancellation hypothesis is hidden
in the proof.

\Cref{prop:valued-amplification} is the case $\delta=D_\sigma$ of
\cref{thm:sigma-derivation}; the case $\sigma=\id$ is used in
\cref{sec:differential} below. The ordered two-probe witness generalizes in
the same way.

\begin{proposition}[Ordered two-point witness]\label[proposition]{prop:sigma-witness}
Let $F$ be an ordered field, $\sigma:F\to F$ a unital field endomorphism,
and $\delta$ a $\sigma$-derivation with $\delta(a)\ne0$. For every $B>0$ put
\[
 b=\frac{2B(1+\abs{\sigma(a)})}{\abs{\delta(a)}}.
\]
Then $\max\{\abs{\delta(b)},\abs{\delta(ab)}\}>B$. Consequently a nonzero
$\sigma$-derivation of an ordered field has cofinal and coinitial image.
\end{proposition}
\begin{proof}
This is the computation proving \cref{lem:amplify}, with
\eqref{eq:sigma-leibniz} in place of \eqref{eq:twisted}. If both absolute
values were at most $B$, then
\[
 2B(1+\abs{\sigma(a)})=\abs{b\delta(a)}
  =\abs{\delta(ab)-\sigma(a)\delta(b)}
  \le B+\abs{\sigma(a)}B,
\]
which is impossible in an ordered field. Since $\delta(-x)=-\delta(x)$,
both signs occur at every attained magnitude, and for a given $c\in F$ one
applies the result with $B>\abs c$.
\end{proof}

Thus \cref{lem:amplify} is the case $\delta=D_\sigma$ and
\cref{lem:derivation-amplify} below is the case $\sigma=\id$. The
unification is not needed for the surreal theorem; it is recorded because
the same two probes settle the displacement and the differential statements
at once.


\section{Faithful action on every nontrivial convex value quotient}

\subsection{Cofinal displacement on the value group}

Suppose now that $\sigma$ is an exponential automorphism stabilizing
$\OO_w$. With $\nu=w\circ E$, there are exact identities
\begin{align}
 \nu(\sigma x)&=\tau_\sigma(\nu(x)),\label{eq:intertwine}\\
 \nu(D_\sigma(x))&=\tau_\sigma(\nu(x))-\nu(x).
 \label{eq:displacement-intertwine}
\end{align}
The first uses $E(\sigma x)=\sigma(E(x))$; the second also uses additivity
of $\nu$. Surjectivity of $\nu$ yields the equality of images
\begin{equation}\label{eq:image-equality}
 \nu(D_\sigma(F))=(\tau_\sigma-\id)(\Gamma).
\end{equation}

\begin{theorem}[Cofinal value displacement]\label[theorem]{thm:cofinal}
Under the hypotheses of \cref{thm:intro}, every nonidentity $\sigma$ has
$(\tau_\sigma-\id)(\Gamma)$ cofinal and coinitial in $\Gamma$.
\end{theorem}
\begin{proof}
By \cref{lem:amplify}, $D_\sigma(F)$ is cofinal and coinitial in $F$.
Fix $\gamma\in\Gamma$ and choose $x\in F$ with $\nu(x)=\gamma$.
There are displacements $d_-<x<d_+$. Order reversal gives
\[
 \nu(d_-)\ge\gamma\ge\nu(d_+).
\]
Thus $\nu(D_\sigma(F))$ is cofinal and coinitial. Equation
\eqref{eq:image-equality} finishes the proof. To obtain strict inequalities
at a given target, first choose a larger, respectively smaller, value in
$\Gamma$, which has no endpoints because it is nontrivial.
\end{proof}

\begin{corollary}[Faithfulness]\label[corollary]{cor:faithful}
The homomorphism
\[
 \rho_w:\Gexp(F;\OO_w)\longrightarrow\Aut(\Gamma,+,<),
 \qquad\sigma\longmapsto\tau_\sigma,
\]
is injective. Here $\Gexp(F;\OO_w)$ denotes the stabilizer of $\OO_w$.
\end{corollary}
\begin{proof}
Compatibility with composition follows by applying $w$ to
$\sigma_1\sigma_2(y)$ and using surjectivity of $w$. If
$\tau_\sigma=\id$, its displacement image is $\{0\}$, which is not cofinal
in a nonzero ordered abelian group. Hence $\sigma=\id$.
\end{proof}

There is also a direct kernel proof. If $\tau_\sigma=\id$,
\eqref{eq:displacement-intertwine} puts every $D_\sigma(x)$ in the proper
convex subgroup $H_w$. Then \cref{cor:bounded-displacement} applies.
The stronger theorem records what happens away from the kernel as well.

\begin{corollary}[Self-embedding version]\label[corollary]{cor:embedding}
Let $\sigma:F\to F$ be a unital field self-embedding commuting with $E$.
If $w(\sigma y)=w(y)$ for all $y\ne0$, then $\sigma=\id$.
\end{corollary}
\begin{proof}
For every $x$, exponential compatibility gives
\[
 \nu(\sigma x-x)=w(\sigma(E(x)))-w(E(x))=0.
\]
Apply \cref{cor:bounded-displacement} to $H_w$.
\end{proof}

This is a self-embedding statement. It does not claim that two embeddings
into an arbitrary larger field coincide merely because they induce the
same source value-group map: the bound and the amplification arguments
must live in the same ambient field as the hypotheses.

\subsection{Rigidity from a single visible exponential scale}

The convention fixed in the introduction --- an ordered exponential is an
increasing group isomorphism $E:(F,+)\to(F_{>0},\cdot)$, so that $E$ is onto
and has an increasing inverse $\log$ --- has been used twice: surjectivity
makes $\nu=w\circ E$ onto $\Gamma$, and the inverse $\log$ appears in
\cref{lem:bridge} and in \cref{prop:probes}. Rigidity itself needs neither.
The following theorem drops surjectivity of $E$, and with it the existence
of $\log$, and replaces nontriviality of $w$ by the assumption that one
positive element has a visible exponential scale.

\begin{theorem}[Rigidity from one visible exponential scale]
\label[theorem]{thm:visible-scale}
Let $F$ be an ordered field, let $w$ be a convex valuation on $F$, and let
\[
 E:(F,+)\longrightarrow(F_{>0},\cdot)
\]
be an order-preserving group homomorphism, \emph{not assumed to be onto}.
Suppose there is one $H>0$ with
\begin{equation}\label{eq:visible-scale}
 w(E(H))<0.
\end{equation}
If a unital field endomorphism $\sigma$ of $F$ satisfies $w(\sigma y)=w(y)$
for every $y\ne0$ and $\sigma(E(x))=E(\sigma x)$ for every $x$, then
$\sigma=\id$.
\end{theorem}
\begin{proof}
Put $u=E(H)>0$. For every $x$, the two hypotheses give
\[
 w\bigl(E(D_\sigma x)\bigr)
  =w\!\left(\frac{E(\sigma x)}{E(x)}\right)
  =w\!\left(\frac{\sigma(E x)}{E x}\right)
  =w(\sigma(Ex))-w(Ex)=0 ,
\]
so $E(D_\sigma x)$ and its inverse $E(-D_\sigma x)$ are positive elements of
$\OO_w$. Suppose $D_\sigma x\ge H$ for some $x$. Then
$0<u\le E(D_\sigma x)$ because $E$ is order-preserving, and convexity in the
form \eqref{eq:convex-monotone} would give
$w(u)\ge w(E(D_\sigma x))=0$, contradicting \eqref{eq:visible-scale}. The
same argument applied to $-D_\sigma x$ rules out $-D_\sigma x\ge H$. Hence
\[
 \abs{D_\sigma x}<H\qquad(x\in F),
\]
and \cref{lem:amplify}, applied with $B=H$, forces $\sigma=\id$.
\end{proof}

The hypothesis removed is surjectivity of $E$. Under the standing
convention the witness \eqref{eq:visible-scale} is automatic whenever $w$ is
nontrivial: choose a positive $u\notin\OO_w$, note that $u>1$, since
$0<u\le1$ would put $u$ in the convex ring, and set $H=\log u>0$, so that
$w(E(H))=w(u)<0$. Consequently \cref{cor:embedding} is the special case of
\cref{thm:visible-scale} in which $E$ is onto and $w$ is nontrivial, and the
theorem is a genuine strengthening: it also covers order-preserving
exponential homomorphisms
whose image is a proper subgroup of $F_{>0}$, for which $\log$ is defined
nowhere outside that image and $\nu$ need not be onto $\Gamma$.

There is no need to identify the convex subgroup $E^{-1}(\OO_w^\times)$ with
anything in particular. For the canonical surreal exponential and the
natural valuation it is the finite surreals, by \cref{prop:finite-log}, but
for a general $E$, or for a coarsening of $w$, it is some other convex
additive subgroup and the proof does not use which one.

\subsection{An order-free, characteristic-free detector}

The same mechanism can be stated with no ordering at all, isolating the
single property of $E$ that rigidity consumes.

\begin{corollary}[Exponential detector]\label[corollary]{cor:detector}
Let $F$ be a field of arbitrary characteristic with a nontrivial surjective
valuation $v:F^\times\to\Gamma$, and let $E:(F,+)\to(F^\times,\cdot)$ be a
group homomorphism. Suppose there is a $\beta\in\Gamma$ such that
\begin{equation}\label{eq:detector}
 E(x)\in\OO_v^\times\quad\Longrightarrow\quad v(x)\ge\beta .
\end{equation}
Then every unital field endomorphism $\sigma$ of $F$ that commutes with $E$
and satisfies $v(\sigma y)=v(y)$ for every $y\ne0$ is the identity.
\end{corollary}
\begin{proof}
For every $x$, the element $E(D_\sigma x)=\sigma(Ex)/Ex$ has value $0$ and
is therefore a unit of $\OO_v$. By \eqref{eq:detector},
$v(D_\sigma x)\ge\beta$ for every $x$. Apply \cref{thm:sigma-derivation}
with $\delta=D_\sigma$, or equivalently \cref{prop:valued-amplification}.
\end{proof}

Neither an ordering, nor convexity of $v$, nor any hypothesis on the
characteristic is used. For the surreal exponential, \cref{prop:finite-log}
supplies \eqref{eq:detector} with $\beta=0$, so \cref{cor:detector} gives a
second proof of \cref{thm:no}, one that never mentions the particular bound
$\omega$. What \eqref{eq:detector} isolates is exactly the property the
rigidity argument consumes: becoming a valuation unit after $E$ must bound
the argument's value from below.


\subsection{Coarsenings and uniqueness of lifts}

For a convex subgroup $\Delta\subseteq\Gamma$, the ordered quotient
$\Gamma/\Delta$ is a value group for the coarsened valuation
\[
 w_\Delta(y)=w(y)+\Delta.
\]
It is nontrivial exactly when $\Delta\ne\Gamma$. The usual correspondence
between convex coarsenings and convex subgroups can also be formulated
for automorphism-invariant valuations; see \cite{DifferenceRank} for this
broader valuation-theoretic setting. The proof below needs only the
quotient formula itself.

\begin{theorem}[Coarsening rigidity]\label[theorem]{thm:coarsening}
Let $\Delta\subsetneq\Gamma$ be convex. If an exponential automorphism
$\sigma$ stabilizes $\OO_w$, stabilizes $\Delta$ through $\tau_\sigma$,
and induces the identity on $\Gamma/\Delta$, then $\sigma=\id$.
Equivalently, the stabilizer of $\Delta$ in $\Gexp(F;\OO_w)$ acts
faithfully on $\Gamma/\Delta$.
\end{theorem}
\begin{proof}
Identity on the quotient says
\[
 (\tau_\sigma-\id)(\Gamma)\subseteq\Delta.
\]
A proper convex subgroup of an ordered abelian group cannot be cofinal:
any positive element outside it is an upper bound. Therefore
\cref{thm:cofinal} excludes $\sigma\ne\id$.
Alternatively, apply \cref{cor:faithful} directly to $w_\Delta$.
\end{proof}

\begin{corollary}[At most one exponential lift]\label[corollary]{cor:unique-lift}
Fix a nontrivial convex valuation, or an invariant nontrivial convex
coarsening, on $F$. An ordered automorphism of its value group has at most
one lift to an exponential automorphism stabilizing that valuation.
\end{corollary}
\begin{proof}
If $\sigma_1$ and $\sigma_2$ induce the same value-group action,
$\sigma_2^{-1}\sigma_1$ induces the identity and hence is the identity.
\end{proof}

The corollary is a uniqueness theorem, not an existence theorem. The next
two sections show that existence is a genuinely restrictive issue.

\subsection{A finite witness at a prescribed scale}

\begin{proposition}[Two exponential probes]\label[proposition]{prop:probes}
Suppose $\sigma(a)\ne a$ and $\eta>0$ belongs to $\Gamma$.
Choose $c>0$ with $w(c)=-\eta$, and set
\[
 B=\log c,\qquad
 b=\frac{2B(1+\abs{\sigma(a)})}{\abs{\sigma(a)-a}}.
\]
Then $B,b>0$, and for at least one $y\in\{b,ab\}$,
\begin{equation}\label{eq:probe}
 \abs{w(\sigma(E(y)))-w(E(y))}\ge\eta.
\end{equation}
The absolute value in \eqref{eq:probe} is the ordered-group absolute value.
\end{proposition}
\begin{proof}
A positive element with negative convex valuation is greater than $1$,
so $B>0$ and $\nu(B)=w(c)=-\eta$.
By \cref{lem:amplify}, one of the two chosen arguments satisfies
$\abs{D_\sigma(y)}>B$.
If $D_\sigma(y)>B$, order reversal gives
$\nu(D_\sigma(y))\le-\eta$; if $D_\sigma(y)<-B$, it gives
$\nu(D_\sigma(y))\ge\eta$.
Use \eqref{eq:displacement-intertwine} and exponential compatibility.
\end{proof}

The result makes the detection mechanism explicit: from one moved element
and a desired scale, two field expressions followed by exponentiation
suffice to exhibit a change at that scale. It is not a finite global
algorithm for deciding equality of arbitrary surreal automorphisms.

\section{The surreal conclusion and the named open question}

\subsection{The complete argument specialized to the surreals}

The following theorem isolates the answer so it can be checked without
using the cofinal-displacement refinement.

\begin{theorem}[No invisible surreal exponential automorphisms]\label[theorem]{thm:no}
Let $\sigma$ be a field automorphism of $\No$ satisfying
\[
 \sigma(\exp x)=\exp(\sigma x)\quad(x\in\No),\qquad
 v(\sigma y)=v(y)\quad(y\ne0).
\]
Then $\sigma=\id$.
\end{theorem}
\begin{proof}
For every $x\in\No$,
\begin{align*}
 v(\exp(\sigma x-x))
 &=v\!\left(\frac{\exp(\sigma x)}{\exp x}\right)\\
 &=v\!\left(\frac{\sigma(\exp x)}{\exp x}\right)=0.
\end{align*}
By \cref{prop:finite-log}, $\sigma x-x$ is finite for every $x$.
In particular,
\[
 \abs{\sigma x-x}<\omega\qquad\text{for every }x\in\No.
\]
If $\sigma$ moved an element $a$, \cref{lem:amplify}, with $B=\omega$,
would produce a displacement with absolute value greater than $\omega$.
This is impossible.
\end{proof}

\begin{corollary}[Answer to Question 5.4]\label[corollary]{cor:question54}
Every exponential $1$-automorphism of $\No$ is the identity.
\end{corollary}
\begin{proof}
Fixing the leading term fixes its exponent, hence fixes the natural
valuation value of every nonzero element. Apply \cref{thm:no}.
\end{proof}

This is the only implication from the definition of a $1$-automorphism that
the proof uses. Leading coefficients play no role. Nor is there an
interchange of $\sigma$ with a Conway-normal-form sum: $\sigma$ is applied
only to individual field expressions.

\subsection{Why the argument is not a claim of total automorphism rigidity}

Every field automorphism of a real-closed field is order preserving:
positive elements are precisely nonzero squares. An ordered field
automorphism also preserves Archimedean comparison, because it preserves
inequalities involving every positive integer. It therefore stabilizes the
natural valuation ring, but it need not fix the value of each element.
Our conclusion says that an exponential automorphism is determined by this
possibly nontrivial action on values.

In particular, it is incorrect to replace the hypothesis
$v(\sigma y)=v(y)$ in \cref{thm:no} by preservation of comparisons
$v(y)<v(z)$. The latter is automatic for ordered field automorphisms. The
faithful action is the precise statement, and it leaves the image of that
action to be understood.

There is another useful conceptual distinction. Suppose one starts with a
field automorphism that changes an element only by an infinitesimal. That
single local observation does not contradict \cref{lem:amplify}.
The contradiction occurs only when a uniform restriction is imposed on the
errors at \emph{all} elements. Multiplication amplifies the error somewhere
else; exponentiation makes the amplified error visible in the value group.

\subsection{A strong nonlocality consequence}

For $\sigma\in\Gexp(\No)$, write $\tau_\sigma$ for its induced automorphism
of the additive group $\No$. If $\sigma\ne\id$, then for every surreal
$M>0$ there are $\gamma_+,\gamma_-\in\No$ such that
\begin{equation}\label{eq:nonlocal}
 \tau_\sigma(\gamma_+)-\gamma_+>M,
 \qquad \tau_\sigma(\gamma_-)-\gamma_-<-M.
\end{equation}
This is \cref{thm:cofinal} with its strict form. The conclusion is stronger
than merely moving some leading exponent. It excludes any value-group
perturbation confined to a bounded range of additive errors, no matter how
large that bound is.

The result does not assert that every individual value has to move, or
that the set of moved values is an interval. It is a statement about the
cofinality of the displacement image. This distinction matters when
constructing and rejecting candidate lifts.

\section{Exponential profiles and hypotheses on a final interval}
\label{sec:profiles}

The hypotheses of \cref{cor:embedding} and \cref{thm:no} are global:
$\sigma$ fixes every value, and $\sigma$ commutes with $E$ at every
argument. The bridge map of \cref{lem:bridge} allows both to be weakened to
a statement about a single final interval. Call
\[
 \nu(x)=w(E(x))
\]
the \emph{exponential profile} of $x$. It is additive with proper convex
kernel $H_w$, and it is very far from injective on individual field
elements. The ambiguity disappears as soon as field multiplication has to
be respected as well.

\begin{theorem}[Profile rigidity]\label[theorem]{thm:profile}
Let $F$ be an ordered field with an ordered exponential $E$ and a nontrivial
convex valuation $w$, and let $\nu=w\circ E$. Let $\sigma:F\to F$ be a
unital field endomorphism. If $S\subseteq F$ generates $(F,+)$ as a group
and
\begin{equation}\label{eq:profile-S}
 \nu(\sigma x)=\nu(x)\qquad(x\in S),
\end{equation}
then $\sigma=\id$. In particular it is enough that \eqref{eq:profile-S}
hold for every $x>c$, where $c\in F$ is one fixed element.
\end{theorem}
\begin{proof}
For $x\in S$, additivity of $\nu$ gives
$\nu(D_\sigma x)=\nu(\sigma x)-\nu(x)=0$, so $D_\sigma x\in H_w$. The
displacement $D_\sigma$ is additive and $H_w$ is an additive subgroup, so
$D_\sigma y\in H_w$ for every $y$ in the subgroup generated by $S$, which is
all of $F$. By \cref{lem:bridge} the subgroup $H_w$ is proper and convex, so
\cref{cor:bounded-displacement} gives $\sigma=\id$.

For the final assertion, take $S=\{x\in F:x>c\}$ and let $x\in F$ be
arbitrary. Put
\[
 z=\abs c+\abs x+1,\qquad y=z+x .
\]
Then $z>c$ because $\abs c\ge c$, and $y\ge\abs c+1>c$ because
$\abs x+x\ge0$. Since $x=y-z$, the final interval generates $(F,+)$.
\end{proof}

\Cref{thm:profile} assumes neither that $\sigma$ fixes the valuation of
field elements nor that it commutes with $E$ anywhere; it compares only the
two profile values in \eqref{eq:profile-S}. The hypotheses of \cref{thm:no}
imply \eqref{eq:profile-S} for every $x$, which is how the general statement
contains the surreal one.

\begin{corollary}[Tail commutation suffices]\label[corollary]{cor:tail}
Let $F$, $E$, $w$ be as in \cref{thm:profile} and let $\sigma$ be a unital
field endomorphism with $w(\sigma y)=w(y)$ for every $y\ne0$. If there is
one $c\in F$ with
\[
 \sigma(E x)=E(\sigma x)\qquad(x>c),
\]
then $\sigma=\id$. The same conclusion follows from the weaker tail
condition $w(\sigma(E x))=w(E(\sigma x))$ for $x>c$.
\end{corollary}
\begin{proof}
Pointwise fixation of $w$ gives $w(\sigma(Ex))=w(Ex)=\nu(x)$. Either tail
condition then yields $\nu(\sigma x)=w(E(\sigma x))=\nu(x)$ for every
$x>c$, so \cref{thm:profile} applies with $S=\{x:x>c\}$.
\end{proof}

\begin{corollary}[Comparing two automorphisms by profile]
\label[corollary]{cor:two-profiles}
Let $\sigma$ and $\tau$ be field automorphisms of $F$ with
$\nu(\sigma x)=\nu(\tau x)$ for every $x\in F$. Then $\sigma=\tau$.
\end{corollary}
\begin{proof}
Additivity of $\nu$ gives $\sigma x-\tau x\in H_w$ for every $x$. Writing
$y=\tau x$ and using surjectivity of $\tau$, we get
$(\sigma\tau^{-1})(y)-y\in H_w$ for every $y\in F$. Now apply
\cref{cor:bounded-displacement}.
\end{proof}

These are global conclusions drawn from coarse measurements, and they must
not be read as saying more. The profile does not reconstruct its argument:
$\nu$ remains highly non-injective, and none of these statements is an
approximation procedure for individual normal forms. Note also that
\cref{cor:two-profiles} uses surjectivity of $\tau$, whereas
\cref{thm:profile}, which compares one endomorphism with the identity,
uses no surjectivity at all.


\section{Multiplicative motion and exponential commutation defects}
\label{sec:defects}

Rigidity says that a nonidentity exponential automorphism cannot be
invisible. The amplification lemma says more: its multiplicative motion,
and its failure to commute with $E$, must appear at arbitrarily large and
arbitrarily small scales.

\subsection{Multiplicative motion}

\begin{theorem}[Multiplicative motion is cofinal]\label[theorem]{thm:motion}
Let $F$ be an ordered field with an ordered exponential $E$, and let
$\sigma\ne\id$ be a unital field endomorphism commuting with $E$. Then
\[
 \left\{\frac{\sigma(y)}{y}:y\in F_{>0}\right\}
\]
is cofinal and coinitial in $F_{>0}$: for every $M>1$ there are positive
$y,z$ with
\[
 \frac{\sigma(y)}{y}>M,\qquad 0<\frac{\sigma(z)}{z}<\frac1M .
\]
\end{theorem}
\begin{proof}
Every $y>0$ equals $E(x)$ for some $x$, because $E$ is onto $F_{>0}$, and
\[
 \frac{\sigma(Ex)}{Ex}=\frac{E(\sigma x)}{Ex}=E(D_\sigma x)>0 .
\]
By \cref{lem:amplify} there are $x$ with $D_\sigma x>\log M$ and $x'$ with
$D_\sigma x'<-\log M$. Applying the increasing map $E$ and putting
$y=E(x)$ and $z=E(x')$ gives the two inequalities.
\end{proof}

No valuation is used here. The statement is that a nonidentity exponential
automorphism cannot be multiplicatively near the identity at every scale.
Surjectivity of $E$ is used, both to represent an arbitrary positive $y$ and
to form $\log M$.

\subsection{The defect when values are fixed but exponentiation is not}

Drop the commutation hypothesis and keep the valuation one. For a unital
field endomorphism $\sigma$ of $F$, define the \emph{exponential commutation
defect}
\begin{equation}\label{eq:defect}
 C_\sigma(x)=\frac{\sigma(E x)}{E(\sigma x)}\qquad(x\in F).
\end{equation}
Numerator and denominator are nonzero, so $w(C_\sigma(x))$ is defined, and
$\sigma$ commutes with $E$ exactly when $C_\sigma$ is identically $1$.

\begin{theorem}[Cofinal commutation defects]\label[theorem]{thm:defect}
Let $F$ be an ordered field with an ordered exponential $E$ and a nontrivial
convex valuation $w$, and let $\sigma\ne\id$ be a unital field endomorphism
with $w(\sigma y)=w(y)$ for every $y\ne0$. Then
\[
 \{w(C_\sigma(x)):x\in F\}
\]
is cofinal and coinitial in $\Gamma$.
\end{theorem}
\begin{proof}
Pointwise fixation of $w$ and additivity of $\nu$ give
\begin{align}
 w(C_\sigma(x))&=w(\sigma(Ex))-w(E(\sigma x))\notag\\
 &=w(Ex)-w(E(\sigma x))
 =\nu(x)-\nu(\sigma x)=-\nu(D_\sigma x).\label{eq:defect-value}
\end{align}
By \cref{lem:amplify} the image $D_\sigma(F)$ is cofinal and coinitial in
$F$, and by \cref{lem:bridge} the map $\nu$ is a surjective order-reversing
additive map. Fix a target $\gamma\in\Gamma$ and choose $\eta>0$ in
$\Gamma$, which is possible since $\Gamma$ is nontrivial. Choose $r\in F$
with $\nu(r)=-\gamma-\eta$ and then $x$ with $D_\sigma x>r$; order reversal
gives $\nu(D_\sigma x)\le-\gamma-\eta$, hence
$w(C_\sigma(x))\ge\gamma+\eta>\gamma$. Choose $s\in F$ with
$\nu(s)=-\gamma+\eta$ and then $x'$ with $D_\sigma x'<s$; order reversal
gives $\nu(D_\sigma x')\ge-\gamma+\eta$, hence
$w(C_\sigma(x'))\le\gamma-\eta<\gamma$.
\end{proof}

\begin{corollary}[No uniformly bounded approximate commutation]
\label[corollary]{cor:approx}
Let $F$, $E$, $w$ be as in \cref{thm:defect} and let $\sigma$ be a unital
field endomorphism with $w(\sigma y)=w(y)$ for every $y\ne0$. If the values
$w(C_\sigma(x))$ are bounded above for all $x$, or bounded below for all
$x$, then $\sigma=\id$.
\end{corollary}
\begin{proof}
Either one-sided bound contradicts \cref{thm:defect}.
\end{proof}

For $F=\No$ this says that a nonidentity exact-value field automorphism
fails to commute with the surreal exponential at unboundedly many
\emph{scales}, in both valuation directions, rather than merely at some one
argument. Two cautions. The word ``approximate'' in \cref{cor:approx}
refers only to the stated valuation bounds; no topology and no numerical
approximation procedure is involved. And \cref{thm:defect} is conditional:
it does not assert that a nonidentity value-fixing field automorphism
exists. By \cref{cor:embedding} no such automorphism commutes with $E$, so
the theorem describes how badly commutation must fail for any that do
exist, not a population known to be nonempty.


\section{Explicit value-group automorphisms with no exponential lift}
\label{sec:nonlifting}

\subsection{Two different normal-form levels}

In this section $\Gamma=(\No,+,<)$ is the value group. An element
$\gamma\in\Gamma$ itself has a Conway normal form
\[
 \gamma=\sum_s r_s\omega^s.
\]
Separately, a field element $x\in\No$ can be written as
\[
 x=\sum_\gamma a_\gamma t^\gamma,
 \qquad t^\gamma=\omega^{-\gamma}.
\]
The index $s$ labels a layer \emph{inside a value}; the index $\gamma$ labels
a monomial \emph{inside a field element}. A reweighting of the first
expansion induces a substitution in the second. Keeping these two levels
separate avoids a common notational confusion.

\subsection{Reweighting a single layer}

Fix $s_0\in\No$ and $c\in\R_{>0}$ with $c\ne1$. Let
$[\gamma]_{s_0}$ denote the real coefficient of $\omega^{s_0}$ in $\gamma$.
Define
\begin{equation}\label{eq:single-layer}
 T_{s_0,c}(\gamma)
  =\gamma+(c-1)[\gamma]_{s_0}\omega^{s_0}.
\end{equation}

\begin{proposition}\label[proposition]{prop:single-layer}
The map $T_{s_0,c}$ is a nonidentity ordered additive automorphism of
$\Gamma$. Its displacement image is exactly $\R\omega^{s_0}$, so it is not
induced by any exponential automorphism of $\No$.
\end{proposition}
\begin{proof}
Coefficient extraction is additive, and the inverse reweights the same
coefficient by $1/c$. The sign of a nonzero surreal is the sign of its
largest-exponent coefficient. That coefficient is either unchanged or
multiplied by the positive real $c$. Thus the map preserves positivity
and hence order.
Since $c-1\ne0$, the coefficient $(c-1)[\gamma]_{s_0}$ ranges over all
real numbers; therefore the displacement image is $\R\omega^{s_0}$.
Every member of this image has absolute value less than
$\omega^{s_0+1}$. It is not cofinal, contrary to \cref{thm:cofinal} for a
nonidentity exponential lift.
\end{proof}

The subgroup $\R\omega^{s_0}$ need not be described as convex. The
argument requires only that it is bounded. For $s_0=0$ and $c=2$, this
example simply doubles the real constant coefficient of each value while
leaving all its other coefficients alone. It is a perfectly valid ordered
group symmetry that exponentiation forbids at the field level.

\subsection{Arbitrary bounded-layer reweightings}

Let $c:\No\to\R_{>0}$ be a specified class function. Define
\begin{equation}\label{eq:weighted}
 T_c\left(\sum_s r_s\omega^s\right)
      =\sum_s c(s)r_s\omega^s.
\end{equation}
Each input has set support, so the output is a legitimate surreal number.
No uniform bound on the positive real factors $c(s)$ is assumed.

\begin{theorem}[Bounded-layer obstruction]\label[theorem]{thm:bounded-layers}
Suppose $c$ is not identically $1$ and there exists $u\in\No$ such that
\[
 c(s)=1\qquad\text{whenever }s\ge u.
\]
Then $T_c$ is an ordered additive automorphism of $\Gamma$, but no
exponential automorphism of $\No$ induces $T_c$ on values.
In particular, this holds whenever the nontrivially reweighted layers
form a nonempty set.
\end{theorem}
\begin{proof}
Additivity, invertibility, and preservation of sign follow as in
\cref{prop:single-layer}; the inverse uses the weights $c(s)^{-1}$.
For any $\gamma$, every exponent in
$\supp(T_c(\gamma)-\gamma)$ is strictly less than $u$. If the difference
is nonzero, its largest exponent is some $s<u$. A real multiple of
$\omega^s$, together with lower terms, has absolute value less than
$\omega^u$. Consequently
\[
 \abs{T_c(\gamma)-\gamma}<\omega^u\qquad(\gamma\in\Gamma).
\]
The displacement is not cofinal, so \cref{thm:cofinal} rules out an
exponential lift.

Every set of surreal numbers has a strict surreal upper bound: for a set
$S$ the Conway cut $\{S\cup\{0\}\mid\varnothing\}$ supplies one.
Thus a set of changed exponents is covered by the stated hypothesis.
\end{proof}

This last clause uses a specifically surreal feature: no set of surreal
exponents is cofinal in $\No$. The bounded-layer theorem itself has the
weaker and more transparent hypothesis of a strict upper bound.

\subsection{The same maps do lift as ordered field automorphisms}

The preceding obstructions are not failures of Hahn substitution. That
substitution works exactly as expected.

\begin{proposition}[Canonical field lift]\label[proposition]{prop:hahn-lift}
Let $T$ be any ordered additive automorphism of $\Gamma=(\No,+,<)$.
The formula
\begin{equation}\label{eq:hahn-lift}
 \widehat T\left(\sum_\gamma a_\gamma t^\gamma\right)
   =\sum_\gamma a_\gamma t^{T(\gamma)}
\end{equation}
defines an ordered field automorphism of $\No$, fixing $\R$ pointwise
and inducing $T$ on its natural value group. It preserves every summable
Hahn family.
\end{proposition}
\begin{proof}
An order isomorphism carries a well-ordered support to a well-ordered
support. The class image of an input set is still a set, so the right-hand
side is a legitimate normal form. Coefficientwise addition is preserved.
For multiplication, the equality
\[
 T(\gamma+\delta)=T(\gamma)+T(\delta)
\]
identifies all contributions to each product coefficient. In a Hahn
product only finitely many support pairs contribute to a fixed exponent;
transport by $T$ gives a bijection between those finite sets of pairs.
Thus multiplication is preserved. The inverse is $\widehat{T^{-1}}$.
The leading coefficient is unchanged and $T$ preserves the ordering of
exponents, so the field order is preserved. Constants have exponent $0$,
which $T$ fixes. Finally, a summable Hahn family has well-ordered union
of supports and only finitely many contributions at any one exponent.
Both properties, and the resulting coefficient sums, are preserved by $T$.
\end{proof}

The construction is the standard external Hahn-field lift; its role in
automorphism decompositions is developed in \cite{KuhlmannSerra}.
For a map in \cref{thm:bounded-layers}, \cref{prop:hahn-lift} is an
explicit field lift, but the non-lifting theorem says more than
$\widehat T\exp\ne\exp\widehat T$: it excludes \emph{every} possible
exponential lift with that value action, including non-strongly additive
ones and lifts that do not fix the embedded real field.

\subsection{A cofinal obstruction: positive rational dilations}

The bounded-layer examples might suggest that cofinal displacement is the
only restriction. It is not.

\begin{theorem}[Rational dilations do not lift]\label[theorem]{thm:dilation}
Let $r\in\Q_{>0}$ with $r\ne1$. The ordered additive automorphism
\[
 T_r:\Gamma\longrightarrow\Gamma,\qquad T_r(\gamma)=r\gamma,
\]
has displacement image equal to all of $\Gamma$ and has the canonical
field lift \eqref{eq:hahn-lift}. Nevertheless, it has no exponential lift
to $\No$.
\end{theorem}
\begin{proof}
Because $r-1$ is a nonzero rational, multiplication by $r-1$ is a bijection
of the divisible group $\Gamma$, so $(T_r-\id)(\Gamma)=\Gamma$.
Assume that an exponential automorphism $\sigma$ induces $T_r$.
The additive map $\nu=v\circ\exp$ is $\Q$-linear: divisibility and
absence of torsion imply $\nu(qx)=q\nu(x)$ for $q\in\Q$.
Hence
\[
 \nu(\sigma x-rx)=r\nu(x)-r\nu(x)=0.
\]
By \cref{prop:finite-log},
\begin{equation}\label{eq:dilation-error}
 \sigma x-rx\in\OO\qquad(x\in\No).
\end{equation}
Choose $x=\omega$ and write
\[
 \sigma x=rx+e,\qquad \sigma(x^2)=rx^2+f,
 \qquad e,f\in\OO.
\]
Multiplicativity gives
\[
 (r^2-r)x^2+2rxe+e^2-f=0.
\]
Dividing by $x^2$ yields
\begin{equation}\label{eq:dilation-contradiction}
 r^2-r+\frac{2re}{x}+\frac{e^2-f}{x^2}=0.
\end{equation}
The last two terms are infinitesimal because $e,f$ are finite and
$x=\omega$. A nonzero rational cannot equal an infinitesimal.
Thus $r^2-r=0$, contrary to $r>0$ and $r\ne1$.
\end{proof}

The rational hypothesis avoids assuming that an arbitrary exponential
automorphism fixes real constants, and avoids any additional assertion
about real linearity of $v\circ\exp$. The proof genuinely excludes all
lifts with the indicated value action, not just the canonical monomial
substitution. It also supplies a multiplication-based obstruction
independent of the bounded-displacement obstruction on the value group.

\section{Surcomplex automorphisms and logarithmic modulus}

\subsection{An intrinsic operation that does not choose a complex logarithm}

Let $F$ now be real closed and equipped with an ordered exponential.
Put $K=F(i)$, where $i^2=-1$. The standard algebraic complexification has
conjugation
\[
 \overline{a+ib}=a-ib
\]
and the $F$-valued modulus
\[
 \abs{a+ib}=\sqrt{a^2+b^2}\in F_{\ge0}.
\]
The nonnegative square root is unique. Its multiplicativity follows by
squaring both sides and using positivity. Define
\begin{equation}\label{eq:L}
 L:K^\times\longrightarrow K,\qquad L(z)=\log\abs z.
\end{equation}
The range is actually contained in $F$. We keep the codomain $K$ when
regarding $L$ as part of a one-sorted algebraic structure; equivalently,
one may work with a relation defining its graph on $K^\times$.

There is no use of $\sin$ or $\cos$ at infinite arguments, no choice of an
argument of $z$, and no complex-logarithm branch. This is compatible with,
but logically separate from, work on surreal exponential subfields and
surcomplex extensions such as \cite{EhrlichKaplan}. Our structure is
specified by \eqref{eq:L}, rather than by an unspecified ``surcomplex
exponential.''

\begin{lemma}[Recovery of the real subfield]\label[lemma]{lem:range}
The image of $L$ is exactly $F$.
\end{lemma}
\begin{proof}
The inclusion $L(K^\times)\subseteq F$ is immediate from the definition.
For every $x\in F$, the positive element $E(x)\in K^\times$ satisfies
$L(E(x))=\log E(x)=x$.
\end{proof}

\subsection{Exact classification of the structured automorphisms}

Write $\GLmod(K)$ for field automorphisms $\Sigma$ satisfying
\[
 \Sigma(L(z))=L(\Sigma(z))\qquad(z\ne0).
\]
No requirement to preserve $F$ or conjugation is imposed in this definition.
Those properties will follow.

\begin{theorem}[Logarithmic-modulus classification]\label[theorem]{thm:L-classification}
Every $\Sigma\in\GLmod(K)$ has the unique form
\begin{equation}\label{eq:Sigma-form}
 \Sigma(a+ib)=\sigma(a)+\varepsilon i\sigma(b),
 \qquad \varepsilon\in\{1,-1\},\quad\sigma\in\Gexp(F).
\end{equation}
Conversely, every map \eqref{eq:Sigma-form} belongs to $\GLmod(K)$.
Consequently,
\[
 \GLmod(K)\simeq\Gexp(F)\times C_2.
\]
\end{theorem}
\begin{proof}
By \cref{lem:range} and commutation with $L$,
\[
 \Sigma(F)=\Sigma(L(K^\times))
           =L(\Sigma(K^\times))=L(K^\times)=F.
\]
Thus $\sigma=\Sigma|_F$ is a field automorphism of $F$.
Because $F$ is real closed, positivity is characterized by nonzero squares,
so $\sigma$ preserves order.
For $y>0$ in $F$, $L(y)=\log y$; consequently
\[
 \sigma(\log y)=\log(\sigma y).
\]
Taking $y=E(x)$ and applying $E$ gives
$\sigma(E(x))=E(\sigma(x))$. Therefore $\sigma\in\Gexp(F)$.
Since $\Sigma(i)^2=-1$ and the two roots of $X^2+1$ in $K$ are $i$ and
$-i$, $\Sigma(i)=\varepsilon i$ for a unique sign. This proves
\eqref{eq:Sigma-form} and uniqueness.

Conversely, let $\sigma\in\Gexp(F)$ and choose $\varepsilon$.
Formula \eqref{eq:Sigma-form} defines a field automorphism by the quadratic
extension construction. Order preservation of $\sigma$ and uniqueness of
positive square roots show that
\[
 \abs{\Sigma(z)}=\sigma(\abs z).
\]
Since $\sigma$ commutes with $\log$ on positive elements,
\[
 L(\Sigma z)=\log(\sigma\abs z)
             =\sigma(\log\abs z)=\Sigma(L(z)).
\]
The coefficientwise extensions of exponential automorphisms commute with
conjugation. Their intersection with $\{\id,\text{conjugation}\}$ is trivial,
which gives the direct rather than a merely semidirect product.
\end{proof}

The theorem recovers the real field from the operation itself. Thus its
surcomplex conclusion is not just a classification obtained by
\emph{assuming in advance} that every automorphism respects a chosen real
axis. The range of $L$ supplies that axis intrinsically. There is a shorter
route to \eqref{eq:Sigma-form} which assumes at the outset that the
automorphism commutes with conjugation and preserves $E$ on the real axis;
it proves strictly less, since \cref{lem:range} derives both of those
properties from $L$ alone, and it is not used here.

\subsection{Extension of the convex valuation}

Let $w$ be a convex valuation on $F$. Put
\begin{equation}\label{eq:wK}
 w_K(z)=w(\abs z)\qquad(z\ne0).
\end{equation}

\begin{lemma}\label[lemma]{lem:wK}
Formula \eqref{eq:wK} is a valuation extending $w$, with the same value
group. For $z=a+ib$,
\begin{equation}\label{eq:min}
 w_K(a+ib)=\min\{w(a),w(b)\},
\end{equation}
where $w(0)=\infty$.
\end{lemma}
\begin{proof}
For $z\ne0$, let $M=\max\{\abs a,\abs b\}>0$.
Ordered-field inequalities give
\[
 M\le\sqrt{a^2+b^2}\le2M.
\]
Since $w(2)=0$, convexity implies
$w(\sqrt{a^2+b^2})=w(M)=\min\{w(a),w(b)\}$.
Multiplicativity of $\abs z$ gives multiplicativity of $w_K$.
For addition, apply the valuation inequality in $F$ to each coordinate
and use \eqref{eq:min}. The restriction to $F$ is $w$, and surjectivity
onto the original value group follows already from this restriction.
\end{proof}

No positive-definite real norm or Archimedean compactness is involved.
The modulus takes values in $F$, and $w_K$ is the algebraic valuation it
induces.

\begin{theorem}[Surcomplex valuation kernel]\label[theorem]{thm:complex-kernel}
Suppose $w$ is nontrivial. In the stabilizer of $\OO_{w_K}$ inside
$\GLmod(K)$, the kernel of the induced action on $\Gamma$ is
\[
 \{\id,\text{conjugation}\}.
\]
More generally, for a proper convex subgroup $\Delta\subsetneq\Gamma$,
the subgroup also stabilizing $\Delta$ has this same kernel on
$\Gamma/\Delta$. Within the subgroup fixing $i$, both actions are faithful.
\end{theorem}
\begin{proof}
By \cref{thm:L-classification}, an automorphism is determined by an
exponential automorphism $\sigma$ of $F$ and a sign. Moreover,
\[
 w_K(\Sigma z)=w(\sigma(\abs z))=\tau_\sigma(w_K(z)),
\]
where stabilization of $\OO_{w_K}$ implies stabilization of
$\OO_w=F\cap\OO_{w_K}$.
Thus the value-group action is exactly that induced by $\sigma$.
If it is the identity, \cref{cor:faithful} gives $\sigma=\id$.
The two choices of sign leave only the identity and conjugation, and both
indeed fix all values. For a coarsening use \cref{thm:coarsening} instead.
Fixing $i$ requires the positive sign.
\end{proof}

\begin{corollary}[Surcomplex application]\label[corollary]{cor:surcomplex}
For $K=\No[i]$ and $L(z)=\log\abs z$, the value-group action of
$\GLmod(K)$ has kernel exactly $\{\id,\text{conjugation}\}$.
The same holds on every invariant nontrivial convex value quotient.
An automorphism preserving $L$ and $i$ is uniquely determined by any such
action.
\end{corollary}

The canonical Hahn lifts from \cref{prop:hahn-lift} also extend
coefficientwise to $\No[i]$ and commute with conjugation. For the forbidden
value-group maps in \cref{thm:bounded-layers,thm:dilation}, these extensions
cannot preserve $L$. More strongly, no $L$-preserving surcomplex
field automorphism has the forbidden value action: restricting it to $\No$
would contradict the corresponding non-lifting theorem.

\section{A differential counterpart}
\label{sec:differential}

The displacement argument also has a useful infinitesimal analogue. This
section is a consequence of the same algebraic mechanism, rather than an
independent priority claim about derivation theory. Formal sums and the
relation between derivations and automorphisms are treated in
\cite{FormalSums}; none of that machinery is needed below.

\begin{lemma}\label[lemma]{lem:derivation-amplify}
If $\partial:F\to F$ is a nonzero derivation of an ordered field, then
$\partial(F)$ is cofinal and coinitial in $F$.
\end{lemma}
\begin{proof}
Choose $a$ with $d=\partial a\ne0$. Given $B>0$, put
\[
 b=\frac{2B(1+\abs a)}{\abs d}.
\]
If $\abs{\partial b}\le B$ and $\abs{\partial(ab)}\le B$, the Leibniz
identity $bd=\partial(ab)-a\partial b$ gives
\[
 2B(1+\abs a)\le B(1+\abs a),
\]
a contradiction. Additivity supplies both signs.
\end{proof}

\begin{theorem}[No globally nonexpanding exponential derivation]\label[theorem]{thm:derivation}
Let $E$ be an ordered exponential on $F$, and let $w$ be a nontrivial convex
valuation. Suppose a derivation $\partial$ satisfies
\[
 \partial E(x)=E(x)\partial x\qquad(x\in F)
\]
and
\begin{equation}\label{eq:nonexpand}
 w(\partial y)\ge w(y)\qquad(y\ne0),
\end{equation}
with the convention $w(0)=\infty$. Then $\partial=0$.
\end{theorem}
\begin{proof}
Apply \eqref{eq:nonexpand} to $y=E(x)$. When $\partial x\ne0$,
\[
 w(E(x))+w(\partial x)=w(\partial E(x))\ge w(E(x)),
\]
so $w(\partial x)\ge0$. The same conclusion holds when $\partial x=0$.
Thus $\partial(F)\subseteq\OO_w$.
A nontrivial valuation has a proper valuation ring, and a convex valuation
ring is a proper convex additive subgroup. It is bounded in the field
order by \cref{lem:bounded}. This contradicts
\cref{lem:derivation-amplify} unless $\partial=0$.
\end{proof}

Condition \eqref{eq:nonexpand} is global. It is much stronger than a
condition imposed only on the valuation ring, such as
$\partial(\OO_w)\subseteq\OO_w$.
The theorem must not be used to rule out surreal exponential derivations
merely because they have a small-derivation property on finite elements.
The obstruction is the universal inequality for \emph{every} nonzero
field element.

\subsection{The same obstruction without an ordering}

\Cref{thm:derivation} forbids the single case $\gamma=0$ of the bound below,
and does so for an ordered field with a convex valuation. Neither the
ordering nor the convexity is needed, and the bound is impossible for every
fixed shift $\gamma$, not only for $\gamma=0$.

\begin{theorem}[Unbounded valuation loss]\label[theorem]{thm:derivation-valued}
Let $F$ be a field with a nontrivial surjective valuation
$v:F^\times\to\Gamma$, let $E:F\to F^\times$ be any map, and let $\partial$
be a derivation of $F$ satisfying
\[
 \partial(E(x))=E(x)\,\partial x\qquad(x\in F).
\]
If there is a single $\gamma\in\Gamma$ with
\begin{equation}\label{eq:loss-bound}
 v(\partial y)\ge v(y)+\gamma\qquad(y\ne0),
\end{equation}
then $\partial=0$. Equivalently, if $\partial\ne0$, then for every
$\gamma\in\Gamma$ there is an $x$ for which $y=E(x)$ satisfies
$v(\partial y)-v(y)<\gamma$.
\end{theorem}
\begin{proof}
Substitute $y=E(x)$ in \eqref{eq:loss-bound} and use compatibility:
\[
 v(E(x))+v(\partial x)=v\bigl(\partial(E(x))\bigr)\ge v(E(x))+\gamma ,
\]
so $v(\partial x)\ge\gamma$ for every $x$, the case $\partial x=0$ being
covered by the convention $v(0)=\infty$. An ordinary derivation is an
$\id$-derivation, so \cref{thm:sigma-derivation} gives $\partial=0$.
For the second form, let $\partial\ne0$; \cref{thm:sigma-derivation}
supplies $x$ with $v(\partial x)<\gamma$, and then $y=E(x)$ satisfies
$v(\partial y)-v(y)=v(\partial x)<\gamma$.
\end{proof}

\begin{corollary}[Rescaling does not help]\label[corollary]{cor:no-contraction}
A nonzero derivation $\partial$ of $F$ satisfying
$\partial(E(x))=E(x)\partial x$ for every $x$ cannot satisfy
$v(\partial y)\ge v(y)$ for every $y\ne0$; in particular it cannot be
globally contracting in the sense that $v(\partial y)>v(y)$ for every
$y\ne0$. The same two statements hold for $c\partial$, for every nonzero
$c\in F$.
\end{corollary}
\begin{proof}
The first assertion is \cref{thm:derivation-valued} with $\gamma=0$. For
$c\ne0$ the map $c\partial$ is again a derivation and satisfies
$(c\partial)(E(x))=E(x)(c\partial)(x)$, so the first assertion applies to it
as well; equivalently, a bound $v((c\partial)y)\ge v(y)$ for all $y\ne0$ is
the bound \eqref{eq:loss-bound} for $\partial$ with $\gamma=-v(c)$.
\end{proof}

The gain in generality is in three places at once: any fixed $\gamma$ rather
than $\gamma=0$, an arbitrary valued field rather than an ordered one with a
convex valuation, and an arbitrary nowhere-vanishing $E$ rather than an
ordered exponential. In particular $E$ is not required to be a group
homomorphism: the proof uses only that $E$ takes nonzero values and
satisfies the displayed compatibility law. What does not change is that the
hypothesis is global. Condition \eqref{eq:loss-bound} compares
$v(\partial y)$ with $v(y)$ at \emph{every} nonzero $y$, and the warning
after \cref{thm:derivation} applies verbatim: a derivation that is merely
small on the finite elements, or contracting on a subring, on a prescribed
support class, or on a bounded domain, is not excluded.
\Cref{sec:laurent} exhibits a nonzero contracting derivation that carries a
compatible \emph{partial} exponential.

\subsection{Application to the surreal derivation}

Berarducci and Mantova construct a nonzero derivation on $\No$ compatible
with the surreal exponential \cite{BM}. \Cref{thm:derivation-valued}
therefore applies to it: for every $\gamma$ there are arguments at which its
valuation gain is smaller than $\gamma$, and by \cref{cor:no-contraction} no
fixed nonzero surreal scalar multiple of it is globally
valuation-contracting. Nothing about the construction is used beyond the
existence of such a derivation and the compatibility law, and neither the
summability properties of that derivation nor its normal-form description is
reproduced here.

This does not contradict the formal operator exponentials of
\cite{FormalSums}. The assertion is that a derivation satisfying both a
total scalar-exponential compatibility law on the whole field and a global
valuation-gain bound must vanish; a derivation used in an operator
construction need not satisfy both of those on all of $\No$.


\section{Sharpness: examples over the formal Laurent series}\label{sec:laurent}

The rigidity mechanism is not a generic property of valued fields, and no
one of its hypotheses can simply be dropped. The examples below are
disjoint from those of \cref{sec:nonlifting}: there the value-group maps
were legitimate and every exponential lift failed, whereas here the field
maps are legitimate and one hypothesis of a rigidity theorem fails.

Throughout this section $L=\R((t))$ is the field of formal Laurent series
over $\R$, ordered so that $t>0$ is infinitesimal, with $t$-adic valuation
$v_t$. The sign of a nonzero element is the sign of its lowest-exponent
coefficient. All series manipulations are coefficientwise and involve only
finitely many contributions at each exponent.

\subsection{Leading-term preservation does not bound displacement}

Define
\begin{equation}\label{eq:substitution}
 \sigma(f(t))=f(t+t^2).
\end{equation}

\begin{proposition}\label[proposition]{prop:substitution}
The map \eqref{eq:substitution} is a nonidentity ordered field automorphism
of $L$ that fixes $\R$ pointwise and fixes the leading term of every nonzero
element, so that $v_t(\sigma f)=v_t(f)$ for all $f\ne0$. For every integer
$m\ge1$,
\begin{equation}\label{eq:laurent-displacement}
 v_t\bigl(\sigma(t^{-m})-t^{-m}\bigr)=1-m .
\end{equation}
\end{proposition}
\begin{proof}
The substitution parameter $p(t)=t+t^2$ has valuation $1$ and leading
coefficient $1$. For a series whose exponents are bounded below, the
coefficient of any fixed power of $t$ in $f(p(t))$ receives contributions
from only finitely many terms of $f$, so the substitution is well defined;
the same finite-coefficient bookkeeping gives preservation of addition and
of multiplication.

For invertibility, put
\[
 h(t)=\frac{\sqrt{1+4t}-1}{2}
 =t-t^2+2t^3-5t^4+14t^5-42t^6+\cdots\in t\R[[t]],
\]
the formal square root being the one with constant coefficient $1$. The
identity $h+h^2=t$ gives $p\circ h=t$. A formal power series with zero
constant coefficient and linear coefficient $1$ has a unique compositional
inverse, its coefficients being determined successively with coefficient
$1$ on the new unknown at each step; hence also $h\circ p=t$, and
substitution by $h$ inverts \eqref{eq:substitution}.

For every integer $n$, $p(t)^n=t^n(1+t)^n$ has leading term $t^n$.
Therefore a leading term $a_nt^n$ of $f$ is carried to itself, which gives
both exact preservation of $v_t$ and preservation of sign, hence of the
order. The map is nonidentity because $\sigma(t)=t+t^2$. Finally,
\[
 \sigma(t^{-m})-t^{-m}=t^{-m}\bigl((1+t)^{-m}-1\bigr)
 =-m\,t^{1-m}+\text{terms of higher valuation},
\]
and $-m\ne0$ in characteristic zero, which is \eqref{eq:laurent-displacement}.
\end{proof}

This automorphism satisfies the \emph{relative} smallness condition
$v_t(\sigma f-f)>v_t(f)$ for every $f\ne0$, with the value $\infty$ allowed
when the difference vanishes. By \eqref{eq:laurent-displacement} its
\emph{absolute} displacement nevertheless has values unbounded below, as
\cref{lem:amplify} and \cref{prop:valued-amplification} require. Relative
leading-term agreement is a much weaker statement than a uniform absolute
bound, and the two must not be exchanged.

\subsection{Commuting with the exponential at every finite argument is not enough}

On $\R[[t]]$ there is the canonical partial exponential
\begin{equation}\label{eq:partial-exp}
 E_0(r+\epsilon)=e^{r}\sum_{n\ge0}\frac{\epsilon^{n}}{n!},
 \qquad r\in\R,\quad\epsilon\in t\R[[t]] .
\end{equation}
The sum is formal: only finitely many powers of $\epsilon$ contribute to
each coefficient. Binomial expansion and finite regrouping in each
coefficient give $E_0(x+y)=E_0(x)E_0(y)$.

\begin{proposition}\label[proposition]{prop:partial-exp}
The nonidentity automorphism \eqref{eq:substitution} commutes with $E_0$ at
every argument in $\R[[t]]$.
\end{proposition}
\begin{proof}
It fixes $r$ and $e^{r}$, and it maps $t\R[[t]]$ into itself, since
$v_t(\epsilon(p(t)))=v_t(\epsilon)$. Each coefficient of either side of
\[
 \sigma\!\left(\sum_{n\ge0}\frac{\epsilon^{n}}{n!}\right)
 =\sum_{n\ge0}\frac{\sigma(\epsilon)^{n}}{n!}
\]
is determined by finitely many indices, and substitution commutes with the
corresponding finite algebra. Multiplying by $e^{r}$ gives the assertion.
\end{proof}

Every value of $E_0$ is a $v_t$-unit, so no visible exponential scale
\eqref{eq:visible-scale} exists inside its domain and the hypothesis of
\cref{thm:visible-scale} fails. This is exactly the local-to-global
inference a rigidity proof must not make: exponential compatibility has to
be applied to the possibly infinite displacement of the amplification
witness, not merely at finite arguments. Under $t=\omega^{-1}$ this Laurent
field is represented inside $\No$ by normal forms; we do not claim that its
substitution automorphism extends to a nonidentity exact-value exponential
automorphism of $\No$, and \cref{thm:no} says that it does not.

\subsection{A purely additive analogue of the amplification lemma is false}

Define an $\R$-linear operator on $L$ by
\[
 N(f)=[t^{1}]f\cdot t^{2},\qquad T=\id+N,
\]
where $[t^{1}]f$ denotes the coefficient of $t$ in $f$.

\begin{proposition}\label[proposition]{prop:additive-bounded}
The map $T$ is a nonidentity order-preserving additive automorphism of $L$
that fixes $1$ and fixes the leading term of every nonzero element, and its
displacement is uniformly bounded:
\[
 \abs{T(f)-f}<t\qquad(f\in L).
\]
It is not multiplicative.
\end{proposition}
\begin{proof}
The coefficient of $t$ in $N(f)$ is zero, so $N^{2}=0$ and
$(\id-N)(\id+N)=\id=(\id+N)(\id-N)$, which gives bijectivity. If
$[t^{1}]f\ne0$, the lowest exponent of $f$ is at most $1$ whereas $N(f)$ has
exponent $2$, so the leading term is unchanged; if $[t^{1}]f=0$, then $T$
fixes $f$. Preservation of the leading term preserves sign, hence order.
For a real coefficient $r$ we have $\abs r t^{2}<t$ because $t$ is
infinitesimal, which is the uniform bound. Finally
\[
 T(t)=t+t^{2},\qquad T(t^{2})=t^{2}\ne(t+t^{2})^{2}=T(t)^{2},
\]
so $T$ is not multiplicative.
\end{proof}

A nonidentity additive automorphism of an ordered field can therefore have
uniformly bounded displacement. \Cref{lem:amplify} and
\cref{cor:bounded-displacement} genuinely consume multiplicativity: the
twisted identity \eqref{eq:twisted} is precisely what $T$ fails to satisfy,
and no additive argument can be substituted for it.

\subsection{A contracting derivation with only a partial exponential}

On the same field put
\begin{equation}\label{eq:laurent-derivation}
 \partial=t^{2}\frac{d}{dt},\qquad
 \partial\Bigl(\sum_{n\ge N}a_nt^{n}\Bigr)=\sum_{n\ge N}n\,a_n t^{n+1}.
\end{equation}
This is a nonzero derivation, and every nonzero $f\in L$ satisfies
\begin{equation}\label{eq:laurent-contraction}
 v_t(\partial f)\ge v_t(f)+1 .
\end{equation}
If the lowest exponent of $f$ is nonzero, its term gives equality; if that
exponent is zero, differentiation removes the term and the valuation
increases further; the case $\partial f=0$ is covered by $v_t(0)=\infty$.
For $x\in t\R[[t]]$, coefficientwise differentiation of \eqref{eq:partial-exp}
gives $\partial(E_0(x))=E_0(x)\partial x$, and this extends to arguments
$c+x$ with $c\in\R$ because $E_0(c+x)=e^{c}E_0(x)$ and $\partial c=0$.
By \cref{thm:derivation-valued} there is nonetheless no map
$E:L\to L^{\times}$ defined on \emph{all} of $L$ with
$\partial(E(x))=E(x)\partial x$: such a map together with
\eqref{eq:laurent-contraction} would force $\partial=0$. The example
pinpoints the global quantifier in that theorem. Formal exponentials of
infinitesimals coexist with contracting derivations; the obstruction
appears only when scalar-exponential compatibility is demanded at every
element of the field.


\section{Foundational scope, assumptions, and proof audit}

\subsection{Set-sized exponential fields}

The main arguments are not dependent on the proper-class size of $\No$.
They apply without change to any set-sized ordered exponential field
carrying a nontrivial convex valuation, and to exponential/logarithmic
subfields of $\No$ when those operations are closed in the subfield.
Work on such subfields includes \cite{EhrlichKaplan,BournezGuilmant}.
An arbitrary fixed Hahn subfield should not be assumed closed under the
surreal exponential or logarithm.

The division of labor among the hypotheses is exact.
The ordered-field structure gives the triangle inequality and positive
scalars used in amplification. The total exponential identifies additive
errors with multiplicative ratios. Its surjectivity makes $\nu$ onto the
whole value group. Convexity makes $\nu$ order-reversing. Nontriviality
prevents the value group from collapsing to zero. Stabilization supplies a
well-defined induced action. No completeness, saturation, differentiability,
or infinite-sum preservation is part of these arguments.
Two of these hypotheses can be traded away. \Cref{thm:visible-scale} replaces
surjectivity of $E$ by a single element $H>0$ with $w(E(H))<0$, and
\cref{cor:detector} drops the ordering and the convexity together, assuming
instead that $E$ turns only arguments of value at least $\beta$ into
valuation units.

\subsection{Class-sized fields and quotients}

For $\No$, an automorphism is itself a class map. Expressions such as
$\Gexp(\No)$ are convenient metamathematical notation for statements about
individual class maps, not assertions that the collection of all such maps
is a set or an NBG class. This is the same convention used in
\cite[Remark~2.1]{KKS}. Every identity and uniqueness statement here can be
read as a statement about specified maps and specified elements.

There is a second, subtler size issue. If $\Delta$ is a proper-class convex
subgroup of the additive group $\No$, its cosets may themselves be proper
classes. One need not form a class whose members are those cosets.
The statement that a map acts identically modulo $\Delta$ means simply
\begin{equation}\label{eq:relational-quotient}
 \tau(\gamma)-\gamma\in\Delta\qquad(\gamma\in\No).
\end{equation}
The proof of \cref{thm:coarsening} uses exactly this relational formulation.
For set-sized value groups the ordinary quotient is literal.

For surreal value groups one can, in addition, realize the quotient by
canonical representatives. Put
\[
 S_\Delta=\{s\in\No:\omega^s\in\Delta\}.
\]
This is an initial segment of the exponent class. Indeed, if $t<s$,
then $0<\omega^t<\omega^s$, so convexity gives membership at $t$.
For a nonzero surreal $\gamma$ with leading exponent $s$, Archimedean
comparison with $\omega^s$ and convexity give
\begin{equation}\label{eq:Delta-support}
 \gamma\in\Delta\quad\Longleftrightarrow\quad s\in S_\Delta.
\end{equation}
For the forward direction, some integer multiple of $\abs\gamma$
exceeds $\omega^s$; for the reverse direction, an integer multiple of
$\omega^s$ exceeds $\abs\gamma$. Both multiples stay in the subgroup.
Thus $\Delta$ consists exactly of the numbers whose support is contained
in $S_\Delta$.

Delete from a normal form all terms whose exponents belong to $S_\Delta$.
The resulting truncation is a surreal, since its support is a subset of a
set support. It is additive, has kernel $\Delta$, and its image is a class
of surreal numbers representing the quotient. Nonzero representatives
have exponents outside $S_\Delta$ and dominate the deleted terms, so their
original ordering gives the quotient ordering. This supplies a concrete
interpretation of $\Gamma/\Delta$ when desired, without making proper
classes into elements.

The same care applies to Hahn lifts: every individual input support is a
set, and its image under a specified class automorphism is a set. No
proper-class sum is formed anywhere in the construction.

\subsection{What has and has not been established}

The core claims are mathematical theorems proved above, not conclusions
from numerical experiments. Their key dependencies can be summarized
without concealing any analytic input:
\begin{center}
\small
\begin{tabular}{@{}p{0.33\textwidth}p{0.58\textwidth}@{}}
\toprule
Result & Essential ingredients\\
\midrule
Displacement amplification & Twisted product identity; ordered-field inequality.\\
Faithful value action & Amplification; the surjective order-reversing map $w\circ E$.\\
Coarsening rigidity & Cofinal displacement; proper convex subgroups are not cofinal.\\
Question 5.4 & $v(\exp x)=0$ exactly for finite $x$; bounded-displacement rigidity.\\
Bounded-layer obstruction & Normal forms; leading-exponent domination; cofinal displacement.\\
Rational-dilation obstruction & Rational linearity of $v\circ\exp$; multiplicativity at $\omega$ and $\omega^2$.\\
Surcomplex kernel & Range of $L$; real closedness; quadratic extension; real-field rigidity.\\
Visible-scale rigidity & One $H>0$ with $w(E(H))<0$; convexity; amplification. No surjectivity of $E$.\\
Twisted-derivation bound & Twisted Leibniz identity; surjectivity of a nontrivial valuation.\\
Profile rigidity & Additivity of $w\circ E$ with proper convex kernel; generation by a final interval.\\
Commutation defects & Amplification; order-reversing surjectivity of $w\circ E$; exact value preservation.\\
Differential loss bound & Twisted-derivation bound at $\sigma=\id$; exponential compatibility of $\partial$.\\
\bottomrule
\end{tabular}
\end{center}

We have not classified the image of $\rho_v$, constructed a global
surreal composition theory, classified omega-map-preserving
automorphisms, or asserted that an arbitrary ordered additive value-group
automorphism lifts exponentially. The non-lifting theorems specifically
show the last assertion to be false.

The accompanying executable check files test finite algebraic identities,
finite Hahn analogues and truncated Laurent-series computations using exact
arithmetic. They provide reproducible protection against elementary formula
and sign mistakes. They are not a model of the full surreal field, they do
not certify the theorems involving all elements or class-sized supports, and
in particular nothing in them bears on the class-sized or transfinite
content of any statement above. No Lean proof is claimed for this
manuscript.

\section{Further questions sharpened by the results}

The faithfulness theorem replaces a kernel problem by an image problem.
The most direct next objective is to characterize the ordered additive
automorphisms $T$ of $\No$ that actually occur as $\tau_\sigma$.
The bounded-layer and rational-dilation examples show that this image is
strictly smaller than the full ordered-group automorphism collection,
and that a cofinal-displacement test alone cannot characterize it.

\begin{question}[Image of the value action]\label[question]{q:image}
Can one describe necessary and sufficient compatibility conditions on
$T\in\Aut(\No,+,<)$ for it to admit an exponential lift?
\end{question}

The uniqueness part is already settled by \cref{cor:unique-lift}.
The obstruction in \cref{thm:dilation} suggests that a useful criterion
must retain multiplicative information transported through
$\nu=v\circ\exp$, rather than only the order structure of the value group.
The quotient $F/H_w$ itself is generally not a field: $H_w$ is an additive
subgroup, not a field ideal. Thus one cannot simply impose ordinary field
multiplication on the quotient and declare the lifting problem solved.

\begin{question}[Effective obstructions for specified substitutions]
For a finitely specified ordered-group substitution $T$, can one find a
finite collection of algebraic/exponential tests that either constructs
its unique exponential lift or certifies non-liftability in a chosen
exponential subfield?
\end{question}

The two-probe result gives a finite witness \emph{once a candidate
nonidentity automorphism is given}; it does not by itself provide the
existence test in this question. The distinction between testing a map
already defined on the field and lifting a prescribed value-group map
must be maintained.

The surcomplex theorem reduces the corresponding image question for
$L$-preserving automorphisms to the real exponential-field question, with
only a sign choice for $i$. This reduction is useful precisely because it
needs no global complex logarithm and no analytic path theory.

\section{Conclusion}

The decisive mechanism is elementary but global. A nonidentity field
endomorphism cannot have its additive errors confined to a bounded
region of its own ordered field. Exponentiation sends those errors to
multiplicative ratios, and a convex valuation reads off their scales.
This makes every nontrivial invariant convex value quotient a faithful
detector of exponential automorphisms.

For $\No$, it follows that there are no nonidentity exponential
$1$-automorphisms, without strong-additivity assumptions. The same argument
also produces cofinal displacement and exact non-lifting criteria that
separate ordered Hahn-field symmetries from exponential symmetries.
For $\No[i]$ equipped with logarithmic modulus, the only symmetry left
invisible by these value quotients is conjugation. These conclusions are
independent of any convergence interpretation of surreal normal forms.

\begin{thebibliography}{99}
\bibitem{KKS}
E.~Kaplan, L.~S.~Krapp, and M.~Serra,
\emph{Decomposing the automorphism group of the surreal numbers},
arXiv:2509.22374v3, version submitted 23 April 2026; PDF dated 27 April 2026.
In particular, Definitions~2.4--2.5, Remark~2.1, Proposition~5.2, and
Question~5.4, p.~10.
\url{https://arxiv.org/abs/2509.22374v3}.

\bibitem{Gonshor}
H.~Gonshor,
\emph{An Introduction to the Theory of Surreal Numbers},
London Mathematical Society Lecture Note Series, vol.~110,
Cambridge University Press, 1986. Chapters~5 and~10.
\href{https://doi.org/10.1017/CBO9780511629143}{doi:10.1017/CBO9780511629143}.

\bibitem{DriesEhrlich}
L.~van den Dries and P.~Ehrlich,
\emph{Fields of surreal numbers and exponentiation},
Fundamenta Mathematicae \textbf{167} (2001), 173--188.
\href{https://doi.org/10.4064/fm167-2-3}{doi:10.4064/fm167-2-3}.
Erratum, \textbf{168} (2001), 295--297.

\bibitem{BM}
A.~Berarducci and V.~Mantova,
\emph{Surreal numbers, derivations and transseries},
Journal of the European Mathematical Society \textbf{20} (2018), 339--390.
\href{https://doi.org/10.4171/JEMS/769}{doi:10.4171/JEMS/769}.
\url{https://arxiv.org/abs/1503.00315v3}.
Cited only for the existence of a nonzero exponential-compatible surreal
derivation.

\bibitem{KuhlmannSerra}
S.~Kuhlmann and M.~Serra,
\emph{The automorphism group of a valued field of generalised formal power series},
Journal of Algebra \textbf{605} (2022), 339--376.
\href{https://doi.org/10.1016/j.jalgebra.2022.04.023}{doi:10.1016/j.jalgebra.2022.04.023}.

\bibitem{DifferenceRank}
S.~Kuhlmann, M.~Matusinski, and F.~Point,
\emph{The valuation difference rank of a quasi-ordered difference field},
arXiv:1301.2611, version~2, 2016.
\url{https://arxiv.org/abs/1301.2611v2}.

\bibitem{EhrlichKaplan}
P.~Ehrlich and E.~Kaplan,
\emph{Surreal ordered exponential fields},
arXiv:2002.07739, first posted 2020; version~3, 2021.
\url{https://arxiv.org/abs/2002.07739}.

\bibitem{BournezGuilmant}
O.~Bournez and Q.~Guilmant,
\emph{Surreal fields stable under exponential and logarithmic functions},
arXiv:2201.08199, 2022.
\url{https://arxiv.org/abs/2201.08199}.

\bibitem{FormalSums}
V.~Bagayoko, L.~S.~Krapp, S.~Kuhlmann, D.~Panazzolo, and M.~Serra,
\emph{Automorphisms and derivations on algebras endowed with formal infinite sums},
arXiv:2403.05827, version~2, 2025.
\url{https://arxiv.org/abs/2403.05827v2}.

\bibitem{Repo}
\begingroup\raggedright
V.~Reshetnikov,
\emph{Surreal}, GitHub repository, documentation map \texttt{docs/README.md}.
Inspected snapshot: commit\\
\texttt{aa846271b4dcae2c055b216126a87210292ec19b}.
\url{https://github.com/VladimirReshetnikov/Surreal}.
The repository is research context, not an external proof source for the
new results of this article.
\par\endgroup
\end{thebibliography}
\end{document}
